%%%%%
%
\section{Rotation 2D vers le rep�re port� par le vent moyen}

\itemstyle{Op�rateur~:} \operfmt{Rotat\_2D\_rep\_vent\_moy}
\index{Rotat\_2D\_rep\_vent\_moy}

%%%%%OBJET%%%%
\itemstyle{Objet~:}
Calcul des 3 composantes du vecteur-vent dans 
un rep�re port� par le vent moyen, qui v�rifie (si U est le module du vent): 
\[
\begin{array}{c}
\overline{u} = U\\
\overline{v} = 0\\
\overline{w} = 0\\
\end{array}
\]
Chaque moyenne est calcul�e sur la p�riode pass�e en argument $T_{int}$ ; 
pour chaque s�quence, on appelle l'op�rateur Rotat\_rep\_vent\_moy avec 
un flag � 0 pour ne faire que les deux premi�res rotations.
Dans la d�termination des flux par la m�thode des covariances, cet 
op�rateur est utilis� avant le detrending (qui enl�ve le signal basse
fr�quence ou la tendance), et avant le calcul des covariances.
 
\bigskip
%%%%%ENTREES%%%%
\begin{tabarg}
\multicolumn{3}{l}{\titretab {Entr\'ees}}\\
\hline
$u$ & Vecteur, \texttt {float} & premi\`ere composante horizontale du vent (any)\\
$v$ & Vecteur, \texttt {float} & seconde composante horizontale du vent (m\^eme unit\'e que $u$)\\
$w$ & Vecteur, \texttt {float} & composante verticale du vent (m\^eme unit\'e que $v$)\\
$F_{e}$ & Coef., \texttt{float} & fr\'equence d'\'echantillonnage (\textit{Hz})\\
$T_{int}$ & Coef.,\texttt{float} & p\'eriode d'int\'egration ou dur\'ee de la s\'equence (\textit{s})\\
\end{tabarg}

\medskip
%%%%%SORTIES%%%%
\begin{tabarg}
\multicolumn{3}{l}{\titretab {Sorties}}\\
\hline
$u_{r}$ & Vecteur, \texttt{float} & premi\`ere composante du vent selon U moyen (m\^eme unit\'e que $u$)\\
$v_{r}$ & Vecteur, \texttt{float} & seconde composante du vent dans le rep\`ere UVW moyen (m\^eme unit\'e que $u$)\\
$w_{r}$ & Vecteur, \texttt{float} & troisi\`eme composante du vent dans le rep\`ere UVW moyen (m\^eme unit\'e que $u$)\\
\end{tabarg}


%%%%%FORMULE%%%%
\itemstyle{Formule}

%%Rotat_rep_vent_moy
Application d'une rotation 2D destin\'ee \`a annuler
les composantes transverses (v) et verticales (w) moyennes du vent ;
  
\[
  \left(
\begin{array}{c}
u_r \\ v_r \\ w_r \\
\end{array}
  \right)
=
  \left(
\begin{array}{ccc}
\cos\epsilon \cos\theta & \sin\epsilon \cos\theta & \sin\theta \\
-\sin\epsilon & \cos\epsilon & 0 \\
-\cos\epsilon \sin\theta & -\sin\epsilon \sin\theta & \cos\theta \\
\end{array}
  \right)
.
  \left(
\begin{array}{c}
u \\ v \\ w \\
\end{array}
  \right)
\]

avec $\cos\epsilon=\frac{\overline{u}}{\sqrt{\overline{u}^2+\overline{v}^2}}$,
    $\sin\epsilon=\frac{\overline{v}}{\sqrt{\overline{u}^2+\overline{v}^2}}$,
    $\cos\theta=\frac{\sqrt{\overline{u}^2+\overline{v}^2}}
                     {\sqrt{\overline{u}^2+\overline{v}^2+\overline{w}^2}}$,
et  $\sin\theta=\frac{\overline{w}}
                     {\sqrt{\overline{u}^2+\overline{v}^2}+\overline{w}^2}$ 
(avec $\theta$ limit\'e \`a 20 degr\'es).



%%%%%COMMENTAIRE%%%%
\itemstyle{ Commentaire~: }

Le rep�re ($u$, $v$, $w$) d'origine peut �tre g�ographique ou li� � l'instrument. 
Il doit �tre orthonorm�.
%%%% Question en suspens : doit-il �tre direct ? pas de raisons...


%%%%%AUTEUR%%%%
\itemstyle{Auteur~:}
CNRM/GMEI/TRAMM

%%%%%REFERENCE%%%%
\itemstyle{R�f�rences~:}
Pages 26 � 28 de la note de P. Bessemoulin sur la
d�termination des flux de surface \cite{GMEI-6}.
